\section{Real-Time Inference Design}
\label{sec:3_11}

Operational deployment requires real-time scoring of new listings as they are submitted. This section details the streaming architecture and API service design.

\subsection{Streaming Architecture}

\textbf{Event Ingestion:} New listing events stream from the marketplace platform via message queue. Each event contains listing attributes and SEON transaction data from the submission-time fraud check.

\textbf{Graph Update Strategy:} Two approaches were evaluated:

\textit{Batch Graph Reconstruction:} Rebuild the full graph periodically (daily/weekly) incorporating new listings. Simple but introduces latency---new listings lack graph context until the next rebuild.

\textit{Incremental Graph Updates:} Add new listing nodes and edges in real-time as events arrive. The internal event store tracks entity-to-listing mappings, enabling immediate edge creation when a new listing shares a device/IP/email with existing listings.

The incremental approach is implemented for production, with batch reconstruction as fallback for consistency validation.

\subsection{API Service Design}

\textbf{FastAPI Microservice:} The inference service exposes REST endpoints for integration with the marketplace platform:

\begin{lstlisting}[language=Python, basicstyle=\ttfamily\fontsize{11}{13.2}\selectfont]
POST /predict
  Input: EventPayload (listing attributes + SEON signals)
  Output: PredictResponse (fraud probability, risk factors)

POST /ingest
  Input: EventPayload
  Output: IngestResponse (confirmation, graph update status)
\end{lstlisting}

\textbf{Internal Event Store:} An in-memory store tracks recent events for graph-like feature computation:
\begin{itemize}
    \item Entity-to-listing mappings (device $\rightarrow$ listings, IP $\rightarrow$ listings, etc.)
    \item Enables \texttt{count\_related(entity\_type, value)} queries.
    \item Supports incremental graph context without full graph access.
\end{itemize}

\textbf{Prediction Flow:}
\begin{enumerate}
    \item Extract tabular features from event payload.
    \item Query event store for graph-like counts (device links, IP links).
    \item If GNN mode: generate embedding using cached GNN model.
    \item Concatenate features, score with XGBoost.
    \item Generate SHAP explanation for risk factors.
    \item Return probability with top contributing features.
\end{enumerate}

\subsection{Latency Considerations}

\textbf{Embedding Precomputation vs Real-Time:} Full GNN inference on the 31M-edge graph is computationally expensive. Two strategies address latency:

\textit{Precomputed Embeddings:} Generate embeddings for all known listings in batch. New listings without precomputed embeddings use neighbor-averaged embeddings or tabular-only scoring.

\textit{Graph-Like Feature Approximation:} Replace relational pattern representations with count-based features (number of device links, IP links) computable in $O(1)$ from the event store. This sacrifices some accuracy for guaranteed low latency.

\textbf{Production Recommendation:} Deploy with graph-like count features for real-time scoring (sub-100ms latency), with full GNN inference in batch for periodic model retraining and historical analysis.

The inference infrastructure enables the pattern-based fraud detection developed in this research to operate in production with acceptable latency while maintaining the ability to leverage graph structure for improved detection.
